\documentclass[UTF8,zihao=-4]{ctexart}
\usepackage[a4paper,margin=1.85cm]{geometry}
\usepackage{hyperref}
\usepackage{verbatim}
\usepackage{xcolor}
\usepackage{fancyhdr}
\hypersetup{hidelinks}
\setlength{\parindent}{0pt}
\setlength{\parskip}{0.45em}
\emergencystretch=4em
\sloppy
\pagestyle{fancy}
\fancyhf{}
\lhead{期末复习 PPT 专项模拟卷}
\rhead{\thepage}
\title{期末复习 PPT 名词解释}
\author{}
\date{}
\begin{document}
\maketitle
\setcounter{tocdepth}{1}
\tableofcontents
\newpage


\section*{使用说明}
\addcontentsline{toc}{section}{使用说明}

本文件用于背诵边角概念和简答题。每个名词按卷面可写的短定义整理。

\section*{一、词法分析}
\addcontentsline{toc}{section}{一、词法分析}

\subsection*{词素}
\addcontentsline{toc}{subsection}{词素}

源程序中实际出现、能够匹配某个词法单元模式的字符序列。例如变量名 \texttt{count} 是一个词素。

\subsection*{词法单元}
\addcontentsline{toc}{subsection}{词法单元}

词法分析器输出的记号类别，例如 \texttt{id}、\texttt{num}、\texttt{if}。一个词法单元可以对应多个不同词素。

\subsection*{正则表达式}
\addcontentsline{toc}{subsection}{正则表达式}

用于描述词法单元模式的形式化表达式。编译器词法分析阶段常用正则表达式描述 token 的结构。

\subsection*{NFA}
\addcontentsline{toc}{subsection}{NFA}

非确定有限自动机。一个状态在同一输入符号上可以有多个后继，也可以有 epsilon 转移。

\subsection*{DFA}
\addcontentsline{toc}{subsection}{DFA}

确定有限自动机。任一状态在任一输入符号上最多只有一个后继状态，执行时只需维护一个当前状态。

\subsection*{epsilon-closure}
\addcontentsline{toc}{subsection}{epsilon-closure}

从某个状态或状态集合出发，只经过 epsilon 边能够到达的所有状态集合。

\subsection*{子集构造法}
\addcontentsline{toc}{subsection}{子集构造法}

把 NFA 的状态集合看成 DFA 的一个状态，从而把 NFA 转换成等价 DFA 的方法。

\subsection*{DFA 最小化}
\addcontentsline{toc}{subsection}{DFA 最小化}

把等价状态合并，使 DFA 状态数减少，同时保持识别语言不变。

\section*{二、自顶向下分析}
\addcontentsline{toc}{section}{二、自顶向下分析}

\subsection*{递归下降分析}
\addcontentsline{toc}{subsection}{递归下降分析}

一种自顶向下语法分析方法。每个非终结符通常对应一个分析过程。

\subsection*{回溯}
\addcontentsline{toc}{subsection}{回溯}

当分析器选择某条产生式失败后，退回之前位置并尝试另一条产生式的过程。

\subsection*{左递归}
\addcontentsline{toc}{subsection}{左递归}

若非终结符 \texttt{A} 能推出以 \texttt{A} 开头的串，则存在左递归。直接左递归形如 \texttt{A -> A alpha}。

\subsection*{左公因子}
\addcontentsline{toc}{subsection}{左公因子}

同一非终结符的多个产生式右部具有相同开头时，这个共同开头称为左公因子。

\subsection*{FIRST 集}
\addcontentsline{toc}{subsection}{FIRST 集}

\texttt{FIRST(alpha)} 是 \texttt{alpha} 能推出的串的第一个终结符集合；若 \texttt{alpha} 能推出空串，则包含 epsilon。

\subsection*{FOLLOW 集}
\addcontentsline{toc}{subsection}{FOLLOW 集}

\texttt{FOLLOW(A)} 是在某些句型中可能紧跟在非终结符 \texttt{A} 后面的终结符集合。

\subsection*{LL(1) 文法}
\addcontentsline{toc}{subsection}{LL(1) 文法}

从左到右扫描输入，构造最左推导，并只向前看一个输入符号即可选择产生式的文法。

\subsection*{预测分析表}
\addcontentsline{toc}{subsection}{预测分析表}

LL(1) 分析使用的表。行是非终结符，列是终结符或 \texttt{\$}，表项填入应使用的产生式。

\section*{三、自底向上分析}
\addcontentsline{toc}{section}{三、自底向上分析}

\subsection*{自底向上分析}
\addcontentsline{toc}{subsection}{自底向上分析}

从输入串开始，不断把产生式右部规约为左部，最终规约到开始符号的分析方法。

\subsection*{句柄}
\addcontentsline{toc}{subsection}{句柄}

在自底向上分析中，当前句型中下一步应该被规约的产生式右部。

\subsection*{最右推导的逆过程}
\addcontentsline{toc}{subsection}{最右推导的逆过程}

自底向上分析的归约序列等价于某个最右推导反过来执行。

\subsection*{LR(0) 项}
\addcontentsline{toc}{subsection}{LR(0) 项}

在产生式右部加入一个点的项目，用于表示当前已经识别到产生式的哪个位置。

\subsection*{规范 LR(0) 项集族}
\addcontentsline{toc}{subsection}{规范 LR(0) 项集族}

由 closure 和 GOTO 构造出的全部 LR(0) 状态集合。

\subsection*{closure}
\addcontentsline{toc}{subsection}{closure}

若项集中某个项目的点后面是非终结符，则把该非终结符的产生式以点在最前的形式加入项集，直到不能继续加入。

\subsection*{GOTO}
\addcontentsline{toc}{subsection}{GOTO}

\texttt{GOTO(I, X)} 表示在项集 \texttt{I} 中读入文法符号 \texttt{X} 后，点越过 \texttt{X} 并做 closure 得到的新项集。

\subsection*{ACTION 表}
\addcontentsline{toc}{subsection}{ACTION 表}

LR 分析表中处理终结符的部分，用于决定移进、规约、接受或报错。

\subsection*{GOTO 表}
\addcontentsline{toc}{subsection}{GOTO 表}

LR 分析表中处理非终结符的部分，用于规约出非终结符后决定转入哪个状态。

\subsection*{移进-规约冲突}
\addcontentsline{toc}{subsection}{移进-规约冲突}

同一状态、同一输入符号下，分析器既可以移进，又可以规约，无法唯一决定动作。

\subsection*{规约-规约冲突}
\addcontentsline{toc}{subsection}{规约-规约冲突}

同一状态、同一输入符号下，分析器可以按两条或多条产生式规约，无法唯一决定动作。

\subsection*{内核项}
\addcontentsline{toc}{subsection}{内核项}

LR(0) 项集中，初始项 \texttt{S' -> . S} 以及所有点不在产生式最左端的项。

\subsection*{非内核项}
\addcontentsline{toc}{subsection}{非内核项}

除初始项 \texttt{S' -> . S} 外，点在产生式最左端的项。

\subsection*{SLR}
\addcontentsline{toc}{subsection}{SLR}

Simple LR。它在 LR(0) 项集族基础上使用 FOLLOW 集限制规约动作。

\subsection*{SLR 的弱点}
\addcontentsline{toc}{subsection}{SLR 的弱点}

FOLLOW 集只与非终结符有关，不区分具体 LR 状态上下文；FOLLOW 集过大时仍可能产生冲突。

\section*{四、语义分析}
\addcontentsline{toc}{section}{四、语义分析}

\subsection*{SDD}
\addcontentsline{toc}{subsection}{SDD}

语法制导定义。它给文法符号关联属性，并用语义规则描述属性如何计算。

\subsection*{SDT}
\addcontentsline{toc}{subsection}{SDT}

语法制导翻译。它把语义动作嵌入产生式，描述分析过程中何时执行动作。

\subsection*{属性}
\addcontentsline{toc}{subsection}{属性}

与文法符号关联的语义信息，例如类型、值、代码、符号表条目等。

\subsection*{综合属性}
\addcontentsline{toc}{subsection}{综合属性}

由子结点或当前结点信息计算，并向父结点传递的属性。

\subsection*{继承属性}
\addcontentsline{toc}{subsection}{继承属性}

由父结点或左侧兄弟结点传给当前结点的属性。

\subsection*{S 属性定义}
\addcontentsline{toc}{subsection}{S 属性定义}

只包含综合属性的 SDD。属性信息只从子结点向父结点传递。

\subsection*{L 属性定义}
\addcontentsline{toc}{subsection}{L 属性定义}

继承属性只能依赖父结点或左侧兄弟结点信息的一类 SDD，适合自左向右求值。

\subsection*{依赖图}
\addcontentsline{toc}{subsection}{依赖图}

表示属性实例之间依赖关系的图。若属性 \texttt{a} 的计算依赖属性 \texttt{b}，则图中有从 \texttt{b} 到 \texttt{a} 的边。

\subsection*{拓扑序}
\addcontentsline{toc}{subsection}{拓扑序}

在无环依赖图中，满足所有依赖先于被依赖对象计算的一种顺序。

\subsection*{属性文法}
\addcontentsline{toc}{subsection}{属性文法}

没有副作用、只通过属性和语义规则描述语义计算的文法扩展。

\subsection*{语义动作}
\addcontentsline{toc}{subsection}{语义动作}

嵌入产生式中的实现动作，例如插入符号表、生成中间代码或输出结果。

\subsection*{副作用}
\addcontentsline{toc}{subsection}{副作用}

语义动作对属性计算之外的外部状态进行修改，例如修改符号表或输出代码。

\subsection*{受控副作用}
\addcontentsline{toc}{subsection}{受控副作用}

执行位置和顺序明确，并且不会破坏属性依赖关系的副作用。

\subsection*{CST}
\addcontentsline{toc}{subsection}{CST}

具体语法树，完整反映文法推导过程，包含全部语法细节。

\subsection*{AST}
\addcontentsline{toc}{subsection}{AST}

抽象语法树，去掉冗余语法细节，只保留表达程序语义所需的核心结构。

\subsection*{AST 结点表示}
\addcontentsline{toc}{subsection}{AST 结点表示}

AST 结点通常用对象或记录表示。内部结点保存 \texttt{op} 和子结点指针，叶子结点保存词法值或标识符名字。

\subsection*{构造 AST 的 SDD}
\addcontentsline{toc}{subsection}{构造 AST 的 SDD}

用语义属性保存子树结点，并在产生式归约或语义动作执行时创建新的 AST 结点。

\subsection*{符号表}
\addcontentsline{toc}{subsection}{符号表}

记录程序中标识符信息的数据结构，通常包含名字、类型、作用域、存储位置等信息。

\subsection*{作用域}
\addcontentsline{toc}{subsection}{作用域}

名字绑定有效的程序区域。内层作用域可以遮蔽外层同名标识符。

\subsection*{符号表栈}
\addcontentsline{toc}{subsection}{符号表栈}

用于处理嵌套作用域的符号表组织方式。进入作用域时压栈，退出作用域时弹栈。

\subsection*{类型检查}
\addcontentsline{toc}{subsection}{类型检查}

根据符号表和表达式类型规则检查赋值、运算、函数调用等是否类型匹配。

\section*{五、中间代码与 IR}
\addcontentsline{toc}{section}{五、中间代码与 IR}

\subsection*{IR}
\addcontentsline{toc}{subsection}{IR}

中间表示。编译器前端和后端之间使用的程序表示形式。

\subsection*{高层 IR}
\addcontentsline{toc}{subsection}{高层 IR}

保留较多源语言结构的 IR，例如循环、函数、对象等结构仍然显式存在。

\subsection*{中层 IR}
\addcontentsline{toc}{subsection}{中层 IR}

介于源语言结构和机器细节之间的 IR。典型形式包括三地址码、四元式、SSA 形式的三地址码。

\subsection*{低层 IR}
\addcontentsline{toc}{subsection}{低层 IR}

接近目标机器的 IR，例如显式跳转、加载/存储、寄存器或栈操作。

\subsection*{树形 IR}
\addcontentsline{toc}{subsection}{树形 IR}

以树结构表示程序，例如 AST。适合前端和早期高层变换，不适合直接做低层优化。

\subsection*{堆栈式 IR}
\addcontentsline{toc}{subsection}{堆栈式 IR}

操作数隐含在栈上的 IR，例如 JVM 字节码。形式紧凑，但数据依赖不如三地址码显式。

\subsection*{强类型 IR}
\addcontentsline{toc}{subsection}{强类型 IR}

在 IR 中保留较精确的类型、对齐或内存访问信息，便于检查和优化。

\subsection*{弱类型 IR}
\addcontentsline{toc}{subsection}{弱类型 IR}

IR 中类型信息较少或较粗略，表示更接近低层机器操作。

\subsection*{CFG}
\addcontentsline{toc}{subsection}{CFG}

控制流图。结点通常是基本块，边表示程序执行时可能发生的控制转移。

\subsection*{DFG}
\addcontentsline{toc}{subsection}{DFG}

数据流图。用于表示值的定义、使用和数据依赖关系。

\subsection*{三地址码}
\addcontentsline{toc}{subsection}{三地址码}

一种常见中间代码形式，每条指令右侧最多包含一个运算符，常见形式为 \texttt{x = y op z}。

\subsection*{四元式}
\addcontentsline{toc}{subsection}{四元式}

三地址码的一种表格表示，字段为 \texttt{(op, arg1, arg2, result)}。

\subsection*{三元式}
\addcontentsline{toc}{subsection}{三元式}

三地址码的一种表示，字段为 \texttt{(op, arg1, arg2)}，结果通过该条指令的位置编号引用。

\subsection*{间接三元式}
\addcontentsline{toc}{subsection}{间接三元式}

在三元式外增加指针表，通过调整指针表控制指令顺序，避免移动三元式本体。

\subsection*{param}
\addcontentsline{toc}{subsection}{param}

函数调用三地址码中用于传递实参的指令。

\subsection*{call}
\addcontentsline{toc}{subsection}{call}

函数调用三地址码中用于执行调用的指令，通常写明函数名和参数个数。

\subsection*{l-value}
\addcontentsline{toc}{subsection}{l-value}

左值，表示可存储位置，通常可以出现在赋值号左侧。

\subsection*{r-value}
\addcontentsline{toc}{subsection}{r-value}

右值，表示表达式的值，通常用于计算或赋值号右侧。

\subsection*{别名}
\addcontentsline{toc}{subsection}{别名}

两个或多个名字或表达式可能引用同一存储位置的情况。

\subsection*{SSA}
\addcontentsline{toc}{subsection}{SSA}

静态单赋值形式。每个变量版本只赋值一次。

\subsection*{phi 函数}
\addcontentsline{toc}{subsection}{phi 函数}

SSA 中用于在控制流汇合处合并不同来源变量版本的特殊记号。

\subsection*{phi 展开}
\addcontentsline{toc}{subsection}{phi 展开}

把 SSA 中的 phi 函数转换成前驱基本块上的复制或移动指令，使后端低层代码不再需要 phi 记号。

\subsection*{短路求值}
\addcontentsline{toc}{subsection}{短路求值}

布尔表达式中根据左侧结果决定是否计算右侧的求值方式。

\subsection*{布尔表达式的控制流用途}
\addcontentsline{toc}{subsection}{布尔表达式的控制流用途}

布尔表达式用于 \texttt{if}、\texttt{while} 等语句条件时，其真假通过跳转到不同代码位置体现。

\subsection*{布尔表达式的逻辑值用途}
\addcontentsline{toc}{subsection}{布尔表达式的逻辑值用途}

布尔表达式出现在赋值右部等位置时，需要计算出明确的逻辑值，例如 \texttt{0} 或 \texttt{1}。

\subsection*{回填}
\addcontentsline{toc}{subsection}{回填}

先生成跳转目标未确定的指令，等目标位置确定后再补上标签或地址。

\subsection*{truelist}
\addcontentsline{toc}{subsection}{truelist}

布尔表达式为真时应跳转、但目标地址尚未确定的指令位置列表。

\subsection*{falselist}
\addcontentsline{toc}{subsection}{falselist}

布尔表达式为假时应跳转、但目标地址尚未确定的指令位置列表。

\subsection*{nextlist}
\addcontentsline{toc}{subsection}{nextlist}

语句执行完后应跳到后继语句、但目标地址尚未确定的指令位置列表。

\subsection*{makelist}
\addcontentsline{toc}{subsection}{makelist}

创建一个只包含某条指令位置的列表，是回填算法使用的辅助函数。

\subsection*{merge}
\addcontentsline{toc}{subsection}{merge}

合并两个回填列表的辅助函数。

\subsection*{nextquad}
\addcontentsline{toc}{subsection}{nextquad}

下一条将要生成的三地址指令编号或位置。

\subsection*{M 标记}
\addcontentsline{toc}{subsection}{M 标记}

回填翻译中的空产生式标记，通常用 \texttt{M.instr = nextquad} 记录后续代码的起始位置。

\subsection*{N 标记}
\addcontentsline{toc}{subsection}{N 标记}

回填翻译中的空产生式标记，通常生成一个 \texttt{goto \_} 指令坯，并把该位置放入 \texttt{N.nextlist}。

\subsection*{穿越}
\addcontentsline{toc}{subsection}{穿越}

在语法制导翻译中通过继承属性向下传递未确定跳转目标，等目标位置确定后再配合回填完成跳转。

\subsection*{fall-through}
\addcontentsline{toc}{subsection}{fall-through}

不显式生成跳转，让控制流自然顺序进入下一条指令。

\subsection*{do-while}
\addcontentsline{toc}{subsection}{do-while}

后测试循环。循环体先执行一次，然后再判断条件。

\subsection*{break}
\addcontentsline{toc}{subsection}{break}

跳出当前循环或 switch 结构，控制流转到结构之后。

\subsection*{continue}
\addcontentsline{toc}{subsection}{continue}

结束当前循环本轮执行，转到下一轮循环的条件检查位置。
\end{document}
